Fix \(\mathfrak E\). By \(\text{def:common-experiment}\), the finite supports, feature map, reference policy, and inverse temperature are public, whereas \(\rho\) remains a coordinate of the law. The first assertion follows from \(\text{lem:fixed-experiment-shell-certificate}\): a law in \(\mathcal M_{\mathfrak E}(d,C,D,\eta)\) exists if and only if at least one of the displayed finitely many exponential nonlinear-semidefinite systems is feasible. This is an equivalence characterization of nonemptiness; no decision algorithm is asserted.

If a nontrivial shell is nonempty, \(\text{thm:localized-upper}\), with the public-experiment index restored, gives a measurable lower-envelope Gibbs learner satisfying
\[
\sup_{P\in\mathcal M_{\mathfrak E}(d,C,D,\eta)}
\mathbb E_{P^{\otimes n}}[J_{\eta,P}(\pi_P^\star)-J_{\eta,P}(\widehat\pi)]
\le\widetilde O\left(\min\left\{\frac{\eta d^2D}{n},\sqrt{\frac{d^2D}{n}}\right\}\right)+n^{-2}.
\]
Since \(\mathfrak R_{n,\mathfrak E}\) is the infimum over all measurable learners, its value is no larger. Inverting the fast and slow branches as in that theorem yields
\[
N^\star_{\mathfrak E}(\epsilon;d,C,D,\eta)
\le\widetilde O\left(d^2D\min\{\eta/\epsilon,\epsilon^{-2}\}\right)
\]
in the stated range. At the boundary, \(\text{thm:zero-risk-boundary}\) gives
\[
N^\star_{\mathfrak E}(\epsilon;d,1,1,\eta)=1
\]
whenever the boundary shell is nonempty.

Now fix \(1<D\le C<e^\eta\). By \(\text{thm:fixed-design-feasibility-and-frontier-nonuniqueness}\), there is a finite public experiment \(\mathfrak E_0\) with a nonempty exact shell and zero minimax risk for every sample size. Hence
\[
N^\star_{\mathfrak E_0}(\epsilon;d,C,D,\eta)=1
\qquad(\epsilon>0).
\]
At the same indices, \(\text{prop:informative-fast-shell-all-indices}\) supplies a finite public experiment \(\mathfrak E_1\) and positive constants \(c(C,D,\eta)\) and \(\epsilon_0(C,D,\eta)\) such that
\[
N^\star_{\mathfrak E_1}(\epsilon;d,C,D,\eta)
\ge c(C,D,\eta)\frac{dD\eta}{\epsilon},
\qquad 0<\epsilon\le\epsilon_0(C,D,\eta).
\]
The constant is unrestricted as a function of \((C,D,\eta)\), so the invariant content of this display is the pointwise local rate \(\Omega_{C,D,\eta}(d/\epsilon)\), not uniform dependence on \(D\) or \(\eta\). The informative proposition assumes only \(1<D\le C<e^\eta\), and therefore also covers \(2D-1\ge e^\eta\). For sufficiently small \(\epsilon\), its lower bound is strictly larger than one, whereas the value under \(\mathfrak E_0\) remains one. Thus \((d,C,D,\eta)\) does not determine fixed-experiment sample complexity.

Finally, let \(L(\epsilon;d,C,D,\eta)\) be any lower bound asserted for every finite public experiment with a nonempty exact shell at the displayed indices. Applying it to \(\mathfrak E_0\) forces
\[
L(\epsilon;d,C,D,\eta)
\le N^\star_{\mathfrak E_0}(\epsilon;d,C,D,\eta)=1.
\]
Therefore no index-only lower bound uniform over public experiments can exceed one or diverge as \(\epsilon\downarrow0\); in particular, a uniform matching \(\Omega(dD\epsilon^{-2})\) branch is false. The argument proves neither a slow branch for a specified informative experiment nor the clean-temperature \(dD\) all-learner upper bound, exactly as stated.